\scanpage{33}
\mspage{31}
% archive sheet number: 32
de deux façons, qui donnent le $=$ résultat \deleted{\unsure{évident|évid.}} (histoire générale de foncteurs adjoints, à expliciter une bonne fois, en \unsure{termes|terme} d'un diagramme de foncteurs
\diagram{\begin{tikzcd}[ampersand replacement=\&]
 X' \arrow[d,leftrightarrow,"{f'^{*}_{*}}"',""{name=vert,right}] \& X \arrow[l,"{q^*}"',""{name=top,below}]\\
 Y' \arrow[from=vert,to=top,bend right=17,shorten >=1pt] \& Y \arrow[l,"{p^*}"'] \arrow[u,leftrightarrow,"{\unsure{F|f}^*}"']
\end{tikzcd}}
en supposant que les foncteurs adj. à g. $f'_*,f_*$ des deux foncteurs verticaux existent).

Un cas intéressant est celui\deleted{\unsure{ci}} où
\[
 pf'=fq,\qquad \unsure{\gamma|\sigma}=\mathrm{id}_{fq}
\]
et où on \deleted{\unsure{peut}} pose $c_\gamma=c$.

\diagram{\begin{tikzcd}[ampersand replacement=\&,column sep=small]
 \text{Ici} \& p^*f_* \arrow[r,"{?}"] \& f'_*q^*\\
 \text{revient à} \& f_* \arrow[r,"{?}"] \& \underbrace{p_*f'_*}_{\deleted{\unsure{f_*|f}}}q^*\simeq\tikz[baseline=(under.base)]{\node[inner sep=0pt] (under) {$(pf')_*$}; \draw ([yshift=1pt]under.south west) -- ([yshift=-1pt]under.south west) -- ([yshift=-1pt]under.south east) -- ([yshift=1pt]under.south east);}q^*\\
 \& \& (fq)_*q^*\simeq(f_*q_*)q^*=f_*(q_*q^*) \arrow[u,"{\gamma_* *q^*}"',"{\unsure{\prime\prime|\shortparallel}}"]\\
 \& \deleted{f_*} \arrow[ur,bend right=10,"{\deleted{f_* *\alpha}}",postaction={decorate,decoration={zigzag,segment length=7pt,amplitude=2pt,post length=15pt}}] \&
\end{tikzcd}}
Or on a $\mathrm{id}\xrightarrow{\alpha}q_*q^*$ (adjonction)

d'où $f_*\xrightarrow{f_* *\alpha}f_*(q_*q^*)$

et en composant avec $\gamma*q^*$, on trouve ce qu'on veut. Et on peut au lieu de ceci chercher
\diagram{\begin{tikzcd}[ampersand replacement=\&,column sep=small]
 f'^*(p^*f_*) \arrow[r,"{?}"] \& q^*\\[-1ex]
 \overset{\unsure{\shortparallel|\prime\prime}}{(pf')^*}f_* \arrow[d,"{\gamma^* * f_*}"] \& \\
 (fq)^*f_*\simeq q^*(f^*\deleted{\unsure{_*}}f_*) \arrow[r,"{q^* *\gamma'}"] \& q^* \& {\text{où }\tikz[baseline=(letter.base)]{\node[inner sep=0pt] (letter) {$\gamma'$}; \node[inner sep=0pt,anchor=north west] (adj) at ([xshift=1pt,yshift=-3pt]letter.south east) {$\scriptstyle\text{adj.}$}; \draw[->] (letter.south west) to[bend right=25] (adj.west);}\colon f^*f_*\longrightarrow\mathrm{id}}
\end{tikzcd}}
